
\documentclass{article}
\usepackage{colortbl}
\usepackage{makecell}
\usepackage{multirow}
\usepackage{supertabular}

\begin{document}

\newcounter{utterance}

\centering \large Interaction Transcript for game `function\_detective', experiment `string\_passive\_examples', episode 163 with gpt{-}5.4.
\vspace{24pt}

{ \footnotesize  \setcounter{utterance}{1}
\setlength{\tabcolsep}{0pt}
\begin{supertabular}{c@{$\;$}|p{.15\linewidth}@{}p{.15\linewidth}p{.15\linewidth}p{.15\linewidth}p{.15\linewidth}p{.15\linewidth}}
    \# & \multicolumn{2}{c}{Player} && \multicolumn{2}{c}{Game Master} \\
    \hline

    \theutterance \stepcounter{utterance}  
    & & & \multicolumn{4}{p{0.6\linewidth}}{
        \cellcolor[rgb]{0.9,0.9,0.9}{
            \makecell[{{p{\linewidth}}}]{
                \texttt{\tiny{[P1$\langle$GM]}}
                \texttt{I have a hidden Python function for you to figure out.} \\
\\ 
\texttt{Your job is to infer what it does and then submit a correct Python implementation.} \\
\\ 
\texttt{Function signature:} \\
\texttt{(text: str) {-}> str} \\
\\ 
\texttt{You have 15 turns total.} \\
\\ 
\texttt{I control which valid input{-}output examples you see. You can request them one at a time.} \\
\\ 
\texttt{You must make your best SOLVE attempt no later than the second{-}to{-}last turn.} \\
\texttt{Do not wait until the very end.} \\
\\ 
\texttt{You can use these actions:} \\
\texttt{{-} JSON format described at the end of the prompt} \\
\texttt{{-} SOLVE: ```python ... ```} \\
\\ 
\texttt{Do not add import statements in your guess.} \\
\texttt{Anything you might need is already available, so just write the function body in hypothesis and return the answer.} \\
\\ 
\texttt{Preloaded examples:} \\
\texttt{1. ('a{-}b', 'a{-}b')} \\
\\ 
\texttt{Scoring:} \\
\texttt{{-} correctness is required for any non{-}zero score} \\
\texttt{{-} in interactive modes, solving with fewer information requests is better} \\
\\ 
\texttt{Input formatting reminders:} \\
\texttt{{-} strings must be quoted} \\
\texttt{{-} booleans must be True/False} \\
\texttt{{-} use commas between multiple values} \\
\texttt{{-} do not write any extra action tags in the same message} \\
\\ 
\texttt{Format example:} \\
\texttt{You write:} \\
\texttt{\{'current\_mode': <probing/confirming/don't know; choose whichever best describes your current state in the task>, 'input\_rationale': <brief description of why you chose this hypothesis>, 'hypothesis': <either the exact string "unknown" if you have no guess yet OR your current best guess as runnable Python matching the given signature, with a one{-}line docstring stating what it does. Never write placeholder code, if you cannot commit to a guess, use "unknown". Match the given signature.} \\
\texttt{e.g.} \\
\texttt{def f(x):} \\
\texttt{	"""Returns the absolute value of x."""} \\
\texttt{	return x if x >= 0 else {-}x} \\
\\ 
\texttt{Use "unknown" if you have no guess.>, 'confidence\_score': <score between 1{-}100. 100 means highest confidence in your hypothesis>\}} \\
\\ 
\texttt{Output only one valid json per turn.} \\
            }
        }
    }
     \\ \\

    \theutterance \stepcounter{utterance}  
    & \multicolumn{4}{p{0.6\linewidth}}{
        \cellcolor[rgb]{0.561,0.737,0.561}{
            \makecell[{{p{\linewidth}}}]{
                \texttt{\tiny{[P1$\rangle$GM]}}
                \texttt{\{"current\_mode":"probing","input\_rationale":"Start by testing whether the function is identity or transforms simple alphabetic text without punctuation.","hypothesis":"unknown","confidence\_score":5\}} \\
\texttt{\{"current\_mode":"probing","input\_rationale":"Test whether the function acts like the identity on a simple string without punctuation.","input":"abc","hypothesis":"unknown","confidence\_score":5\}} \\
            }
        }
    }
    & & \\ \\

    \theutterance \stepcounter{utterance}  
    & & & \multicolumn{4}{p{0.6\linewidth}}{
        \cellcolor[rgb]{0.9,0.9,0.9}{
            \makecell[{{p{\linewidth}}}]{
                \texttt{\tiny{[P1$\langle$GM]}}
                \texttt{Game Violation: ParseError: Response is not valid JSON.} \\
\texttt{Output ONLY:} \\
\texttt{One JSON with confidence\_score, current\_mode, input\_rationale, and hypothesis keys ...} \\
\texttt{OR} \\
\texttt{SOLVE: ```python ... ```} \\
            }
        }
    }
     \\ \\

    \theutterance \stepcounter{utterance}  
    & \multicolumn{4}{p{0.6\linewidth}}{
        \cellcolor[rgb]{0.561,0.737,0.561}{
            \makecell[{{p{\linewidth}}}]{
                \texttt{\tiny{[P1$\rangle$GM]}}
                \texttt{\{"current\_mode":"probing","input\_rationale":"Please provide the output for input 'abc' to test whether the function is the identity on plain alphabetic text.","hypothesis":"unknown","confidence\_score":5\}} \\
            }
        }
    }
    & & \\ \\

    \theutterance \stepcounter{utterance}  
    & & & \multicolumn{4}{p{0.6\linewidth}}{
        \cellcolor[rgb]{0.9,0.9,0.9}{
            \makecell[{{p{\linewidth}}}]{
                \texttt{\tiny{[P1$\langle$GM]}}
                \texttt{OUTPUT: ('hello', 'llo')} \\
            }
        }
    }
     \\ \\

    \theutterance \stepcounter{utterance}  
    & \multicolumn{4}{p{0.6\linewidth}}{
        \cellcolor[rgb]{0.561,0.737,0.561}{
            \makecell[{{p{\linewidth}}}]{
                \texttt{\tiny{[P1$\rangle$GM]}}
                \texttt{\{"current\_mode":"probing","input\_rationale":"The examples suggest the function may drop leading characters in some cases but preserve punctuation{-}containing strings. Please provide the output for input 'hi'.","hypothesis":"def f(text):\textbackslash{}n    \textbackslash{}"\textbackslash{}"\textbackslash{}"Possibly removes some leading characters depending on the text.\textbackslash{}"\textbackslash{}"\textbackslash{}"\textbackslash{}n    return text","confidence\_score":13\}} \\
            }
        }
    }
    & & \\ \\

    \theutterance \stepcounter{utterance}  
    & & & \multicolumn{4}{p{0.6\linewidth}}{
        \cellcolor[rgb]{0.9,0.9,0.9}{
            \makecell[{{p{\linewidth}}}]{
                \texttt{\tiny{[P1$\langle$GM]}}
                \texttt{OUTPUT: ('   ', '   ')} \\
            }
        }
    }
     \\ \\

    \theutterance \stepcounter{utterance}  
    & \multicolumn{4}{p{0.6\linewidth}}{
        \cellcolor[rgb]{0.561,0.737,0.561}{
            \makecell[{{p{\linewidth}}}]{
                \texttt{\tiny{[P1$\rangle$GM]}}
                \texttt{\{"current\_mode":"probing","input\_rationale":"Current evidence fits a rule that may remove the first two characters only for certain strings, perhaps lowercase letters without punctuation or spaces. Please provide the output for input 'he'.","hypothesis":"def f(text):\textbackslash{}n    \textbackslash{}"\textbackslash{}"\textbackslash{}"May remove the first two characters for some alphabetic strings while leaving others unchanged.\textbackslash{}"\textbackslash{}"\textbackslash{}"\textbackslash{}n    return text{[}2:{]} if False else text","confidence\_score":18\}} \\
            }
        }
    }
    & & \\ \\

    \theutterance \stepcounter{utterance}  
    & & & \multicolumn{4}{p{0.6\linewidth}}{
        \cellcolor[rgb]{0.9,0.9,0.9}{
            \makecell[{{p{\linewidth}}}]{
                \texttt{\tiny{[P1$\langle$GM]}}
                \texttt{OUTPUT: ('aba', 'aba')} \\
            }
        }
    }
     \\ \\

    \theutterance \stepcounter{utterance}  
    & \multicolumn{4}{p{0.6\linewidth}}{
        \cellcolor[rgb]{0.561,0.737,0.561}{
            \makecell[{{p{\linewidth}}}]{
                \texttt{\tiny{[P1$\rangle$GM]}}
                \texttt{\{"current\_mode":"confirming","input\_rationale":"A consistent hypothesis is: return the substring starting at the first position whose character appears again later in the string; if no character repeats later, return the original string. Please provide the output for input 'abcb'.","hypothesis":"def f(text):\textbackslash{}n    \textbackslash{}"\textbackslash{}"\textbackslash{}"Returns the suffix beginning at the first character that appears again later; if none do, returns the original string.\textbackslash{}"\textbackslash{}"\textbackslash{}"\textbackslash{}n    for i, ch in enumerate(text):\textbackslash{}n        if ch in text{[}i+1:{]}:\textbackslash{}n            return text{[}i:{]}\textbackslash{}n    return text","confidence\_score":72\}} \\
            }
        }
    }
    & & \\ \\

    \theutterance \stepcounter{utterance}  
    & & & \multicolumn{4}{p{0.6\linewidth}}{
        \cellcolor[rgb]{0.9,0.9,0.9}{
            \makecell[{{p{\linewidth}}}]{
                \texttt{\tiny{[P1$\langle$GM]}}
                \texttt{OUTPUT: ('ab', 'ab')} \\
            }
        }
    }
     \\ \\

    \theutterance \stepcounter{utterance}  
    & \multicolumn{4}{p{0.6\linewidth}}{
        \cellcolor[rgb]{0.561,0.737,0.561}{
            \makecell[{{p{\linewidth}}}]{
                \texttt{\tiny{[P1$\rangle$GM]}}
                \texttt{\{"current\_mode":"probing","input\_rationale":"The strongest hypothesis now is that it returns the suffix starting at the first non{-}space character that is immediately repeated (e.g. 'hello' {-}> 'llo' because of 'll'), and otherwise returns the original string. Please provide the output for input 'book'.","hypothesis":"def f(text):\textbackslash{}n    \textbackslash{}"\textbackslash{}"\textbackslash{}"Returns the suffix starting at the first non{-}space character that is immediately repeated; otherwise returns the original string.\textbackslash{}"\textbackslash{}"\textbackslash{}"\textbackslash{}n    for i in range(len(text) {-} 1):\textbackslash{}n        if text{[}i{]} == text{[}i + 1{]} and text{[}i{]} != ' ':\textbackslash{}n            return text{[}i:{]}\textbackslash{}n    return text","confidence\_score":78\}} \\
            }
        }
    }
    & & \\ \\

    \theutterance \stepcounter{utterance}  
    & & & \multicolumn{4}{p{0.6\linewidth}}{
        \cellcolor[rgb]{0.9,0.9,0.9}{
            \makecell[{{p{\linewidth}}}]{
                \texttt{\tiny{[P1$\langle$GM]}}
                \texttt{OUTPUT: (' ', ' ')} \\
            }
        }
    }
     \\ \\

    \theutterance \stepcounter{utterance}  
    & \multicolumn{4}{p{0.6\linewidth}}{
        \cellcolor[rgb]{0.561,0.737,0.561}{
            \makecell[{{p{\linewidth}}}]{
                \texttt{\tiny{[P1$\rangle$GM]}}
                \texttt{\{"current\_mode":"confirming","input\_rationale":"Please provide another example that can distinguish whether only repeated non{-}space adjacent characters trigger trimming, such as a word with 'oo' in the middle.","hypothesis":"def f(text):\textbackslash{}n    \textbackslash{}"\textbackslash{}"\textbackslash{}"Returns the suffix starting at the first non{-}space character that is immediately repeated; otherwise returns the original string.\textbackslash{}"\textbackslash{}"\textbackslash{}"\textbackslash{}n    for i in range(len(text) {-} 1):\textbackslash{}n        if text{[}i{]} == text{[}i + 1{]} and text{[}i{]} != ' ':\textbackslash{}n            return text{[}i:{]}\textbackslash{}n    return text","confidence\_score":81\}} \\
            }
        }
    }
    & & \\ \\

    \theutterance \stepcounter{utterance}  
    & & & \multicolumn{4}{p{0.6\linewidth}}{
        \cellcolor[rgb]{0.9,0.9,0.9}{
            \makecell[{{p{\linewidth}}}]{
                \texttt{\tiny{[P1$\langle$GM]}}
                \texttt{OUTPUT: ('hello world', 'rld')} \\
            }
        }
    }
     \\ \\

    \theutterance \stepcounter{utterance}  
    & \multicolumn{4}{p{0.6\linewidth}}{
        \cellcolor[rgb]{0.561,0.737,0.561}{
            \makecell[{{p{\linewidth}}}]{
                \texttt{\tiny{[P1$\rangle$GM]}}
                \texttt{\{"current\_mode":"probing","input\_rationale":"A plausible rule is: examine the last whitespace{-}separated word; if it is alphabetic and starts with one or more consonants followed by a vowel, return the part after that first vowel, otherwise return the original string. Please provide the output for input 'a hello'.","hypothesis":"def f(text):\textbackslash{}n    \textbackslash{}"\textbackslash{}"\textbackslash{}"Returns the tail of the last word after its initial consonant cluster and first vowel; otherwise returns the original string.\textbackslash{}"\textbackslash{}"\textbackslash{}"\textbackslash{}n    parts = text.split()\textbackslash{}n    if not parts:\textbackslash{}n        return text\textbackslash{}n    word = parts{[}{-}1{]}\textbackslash{}n    if not word.isalpha() or not word:\textbackslash{}n        return text\textbackslash{}n    vowels = 'aeiouAEIOU'\textbackslash{}n    if word{[}0{]} in vowels:\textbackslash{}n        return text\textbackslash{}n    i = 0\textbackslash{}n    while i < len(word) and word{[}i{]} not in vowels:\textbackslash{}n        i += 1\textbackslash{}n    if 0 < i < len(word):\textbackslash{}n        return word{[}i+1:{]}\textbackslash{}n    return text","confidence\_score":46\}} \\
            }
        }
    }
    & & \\ \\

    \theutterance \stepcounter{utterance}  
    & & & \multicolumn{4}{p{0.6\linewidth}}{
        \cellcolor[rgb]{0.9,0.9,0.9}{
            \makecell[{{p{\linewidth}}}]{
                \texttt{\tiny{[P1$\langle$GM]}}
                \texttt{OUTPUT: ('A', 'A')} \\
            }
        }
    }
     \\ \\

    \theutterance \stepcounter{utterance}  
    & \multicolumn{4}{p{0.6\linewidth}}{
        \cellcolor[rgb]{0.561,0.737,0.561}{
            \makecell[{{p{\linewidth}}}]{
                \texttt{\tiny{[P1$\rangle$GM]}}
                \texttt{\{"current\_mode":"confirming","input\_rationale":"My best current hypothesis is: look at the last whitespace{-}separated word; if it is lowercase alphabetic and begins with consonant+vowel, return that word without its first two characters, otherwise return the original string. A useful distinguishing example would involve a last word starting with two consonants, like 'street' or 'glow'.","hypothesis":"def f(text):\textbackslash{}n    \textbackslash{}"\textbackslash{}"\textbackslash{}"Returns the last word without its first two characters when that last word is lowercase alphabetic and starts with consonant{-}vowel; otherwise returns the original string.\textbackslash{}"\textbackslash{}"\textbackslash{}"\textbackslash{}n    parts = text.split()\textbackslash{}n    if not parts:\textbackslash{}n        return text\textbackslash{}n    word = parts{[}{-}1{]}\textbackslash{}n    vowels = 'aeiou'\textbackslash{}n    if len(word) >= 2 and word.isalpha() and word.islower() and word{[}0{]} not in vowels and word{[}1{]} in vowels:\textbackslash{}n        return word{[}2:{]}\textbackslash{}n    return text","confidence\_score":58\}} \\
            }
        }
    }
    & & \\ \\

    \theutterance \stepcounter{utterance}  
    & & & \multicolumn{4}{p{0.6\linewidth}}{
        \cellcolor[rgb]{0.9,0.9,0.9}{
            \makecell[{{p{\linewidth}}}]{
                \texttt{\tiny{[P1$\langle$GM]}}
                \texttt{OUTPUT: ('aA', 'aA')} \\
            }
        }
    }
     \\ \\

    \theutterance \stepcounter{utterance}  
    & \multicolumn{4}{p{0.6\linewidth}}{
        \cellcolor[rgb]{0.561,0.737,0.561}{
            \makecell[{{p{\linewidth}}}]{
                \texttt{\tiny{[P1$\rangle$GM]}}
                \texttt{\{"current\_mode":"probing","input\_rationale":"Key uncertainty: whether the function returns the last 3 characters of the last lowercase 5{-}letter word, or more generally removes the first 2 characters from the last lowercase alphabetic word of length at least 5. Please provide the output for input 'abcdef'.","hypothesis":"def f(text):\textbackslash{}n    \textbackslash{}"\textbackslash{}"\textbackslash{}"Possibly returns a suffix of the last lowercase alphabetic word, with known cases 'hello'{-}>'llo' and 'hello world'{-}>'rld'.\textbackslash{}"\textbackslash{}"\textbackslash{}"\textbackslash{}n    parts = text.split()\textbackslash{}n    if not parts:\textbackslash{}n        return text\textbackslash{}n    word = parts{[}{-}1{]}\textbackslash{}n    if word.isalpha() and word.islower() and len(word) >= 5:\textbackslash{}n        return word{[}2:{]}\textbackslash{}n    return text","confidence\_score":41\}} \\
            }
        }
    }
    & & \\ \\

    \theutterance \stepcounter{utterance}  
    & & & \multicolumn{4}{p{0.6\linewidth}}{
        \cellcolor[rgb]{0.9,0.9,0.9}{
            \makecell[{{p{\linewidth}}}]{
                \texttt{\tiny{[P1$\langle$GM]}}
                \texttt{OUTPUT: ('a', 'a')} \\
            }
        }
    }
     \\ \\

    \theutterance \stepcounter{utterance}  
    & \multicolumn{4}{p{0.6\linewidth}}{
        \cellcolor[rgb]{0.561,0.737,0.561}{
            \makecell[{{p{\linewidth}}}]{
                \texttt{\tiny{[P1$\rangle$GM]}}
                \texttt{\{"current\_mode":"confirming","input\_rationale":"Best{-}fit rule from all seen examples: take the last whitespace{-}separated word; if it consists only of lowercase letters and its first two letters are consonant then vowel, return that word with those first two letters removed; otherwise return the original text. One more useful confirming example would be a lowercase word like 'cat' or 'there'.","hypothesis":"def f(text):\textbackslash{}n    \textbackslash{}"\textbackslash{}"\textbackslash{}"Returns the last word without its first two letters when that last word is lowercase alphabetic and starts with consonant{-}vowel; otherwise returns the original text.\textbackslash{}"\textbackslash{}"\textbackslash{}"\textbackslash{}n    parts = text.split()\textbackslash{}n    if not parts:\textbackslash{}n        return text\textbackslash{}n    word = parts{[}{-}1{]}\textbackslash{}n    vowels = 'aeiou'\textbackslash{}n    if len(word) >= 2 and word.isalpha() and word.islower() and word{[}0{]} not in vowels and word{[}1{]} in vowels:\textbackslash{}n        return word{[}2:{]}\textbackslash{}n    return text","confidence\_score":67\}} \\
            }
        }
    }
    & & \\ \\

    \theutterance \stepcounter{utterance}  
    & & & \multicolumn{4}{p{0.6\linewidth}}{
        \cellcolor[rgb]{0.9,0.9,0.9}{
            \makecell[{{p{\linewidth}}}]{
                \texttt{\tiny{[P1$\langle$GM]}}
                \texttt{OUTPUT: ('123', '123')} \\
            }
        }
    }
     \\ \\

    \theutterance \stepcounter{utterance}  
    & \multicolumn{4}{p{0.6\linewidth}}{
        \cellcolor[rgb]{0.561,0.737,0.561}{
            \makecell[{{p{\linewidth}}}]{
                \texttt{\tiny{[P1$\rangle$GM]}}
                \texttt{SOLVE: ```python} \\
\texttt{def f(text):} \\
\texttt{    """Returns the last three characters of the input string, or the whole string if shorter."""} \\
\texttt{    return text{[}{-}3:{]}} \\
\texttt{```} \\
            }
        }
    }
    & & \\ \\

    \theutterance \stepcounter{utterance}  
    & & & \multicolumn{2}{p{0.3\linewidth}}{
        \cellcolor[rgb]{0.9,0.9,0.9}{
            \makecell[{{p{\linewidth}}}]{
                \texttt{\tiny{[GM$|$GM]}}
                \texttt{Game Verdict: WIN} \\
            }
        }
    }
    & & \\ \\

    \theutterance \stepcounter{utterance}  
    & & & \multicolumn{2}{p{0.3\linewidth}}{
        \cellcolor[rgb]{0.9,0.9,0.9}{
            \makecell[{{p{\linewidth}}}]{
                \texttt{\tiny{[GM$|$GM]}}
                \texttt{That is correct!} \\
            }
        }
    }
    & & \\ \\

    \theutterance \stepcounter{utterance}  
    & & & \multicolumn{2}{p{0.3\linewidth}}{
        \cellcolor[rgb]{0.9,0.9,0.9}{
            \makecell[{{p{\linewidth}}}]{
                \texttt{\tiny{[GM$|$GM]}}
                \texttt{def take\_last\_three(text: str) {-}> str:} \\
\texttt{    return text{[}{-}3:{]}} \\
            }
        }
    }
    & & \\ \\

    \theutterance \stepcounter{utterance}  
    & & & \multicolumn{2}{p{0.3\linewidth}}{
        \cellcolor[rgb]{0.9,0.9,0.9}{
            \makecell[{{p{\linewidth}}}]{
                \texttt{\tiny{[GM$|$GM]}}
                \texttt{def f(text):} \\
\texttt{    """Returns the last three characters of the input string, or the whole string if shorter."""} \\
\texttt{    return text{[}{-}3:{]}} \\
            }
        }
    }
    & & \\ \\

    \theutterance \stepcounter{utterance}  
    & & & \multicolumn{2}{p{0.3\linewidth}}{
        \cellcolor[rgb]{0.9,0.9,0.9}{
            \makecell[{{p{\linewidth}}}]{
                \texttt{\tiny{[GM$|$GM]}}
                \texttt{Mode: passive\_examples} \\
\texttt{Accuracy: 1.000} \\
\texttt{Information requests: 10/15} \\
\texttt{Revealed passive examples: 11} \\
\texttt{Efficiency: 0.400} \\
\texttt{Slack: 5} \\
\texttt{Observed{-}pair consistency: 1.000} \\
\texttt{Parse errors: 1} \\
\texttt{Runtime errors: 0} \\
            }
        }
    }
    & & \\ \\

\end{supertabular}
}

\end{document}
